% context: 9th-grade physics do-now practice, Unit 2 kinematics, after 2-5 free fall
% topic: mental math, no calculator, answers only — metric length prefixes, squares
%        and square roots, m/s and km/h, distance = speed x time, average speed vs
%        average velocity, fractions to decimals, fractions of an hour,
%        acceleration = change in velocity / time, free fall with g = 10 m/s^2
% structure: two pages; page 1 prefixes, squares, speed units, and speed-time-distance;
%        page 2 fractions, time, acceleration, and falling. Follows 2-1DN and 2-5HW with
%        new numbers, a little easier
% data: standard 400 m running track; marathon 42.195 km (contexts.md #4)
\documentclass[12pt, twoside]{article}
%\documentclass[12pt, twoside]{article}
\usepackage[letterpaper, margin=1in, headsep=0.2in]{geometry}
\setlength{\headheight}{0.6in}
%\usepackage[english]{babel}
\usepackage[utf8]{inputenc}
\usepackage{microtype}
\usepackage{amsmath}
\usepackage{amssymb}
%\usepackage{amsfonts}
\usepackage[nomessages]{fp} %\FPeval{\var-name}{2*sin(pi/6)}
\usepackage{siunitx} %units in math. eg 20\milli\meter
\usepackage{yhmath} % for arcs, overparenth command
\usepackage{tikz} %graphics
\usetikzlibrary{quotes, angles, arrows, arrows.meta}
\usepackage{graphicx} %consider setting \graphicspath{{images/}}
\usepackage{parskip} %no paragraph indent
\usepackage{enumitem}
\usepackage{multicol}
\usepackage{venndiagram}

\usepackage{fancyhdr}
\pagestyle{fancy}
\fancyhf{}
\renewcommand{\headrulewidth}{0pt} % disable the underline of the header
\raggedbottom
\hfuzz=2mm %suppresses overfull box warnings

\usepackage{hyperref}
\usepackage{float}

\newcommand{\blank}[1][2.4cm]{\rule{#1}{0.4pt}}

\title{Physics}
\author{Chris Huson}
\date{November 2026}

\fancyhead[LE]{\thepage}
\fancyhead[RO]{\thepage \\ Nickname: \hspace{2.25cm} \,\\ \,}
\fancyhead[LO]{La Scuola d'Italia / Huson / Physics \\* 5 November 2026}

\begin{document}

\subsubsection*{2.6 Warm-up: Mental Math (no calculator)}

Work each question in your head --- no calculator, and no need to show work. Write
just the answer, with its units.

\vspace{0.15cm}

\textbf{Reference.} \quad
kilo (k) = 1000 \quad centi (c) = $\dfrac{1}{100}$ \quad
milli (m) = $\dfrac{1}{1000}$ \quad 1 m/s = 3.6 km/h

\vspace{0.2cm}

\begin{enumerate}[itemsep=0.35cm]

\item Fill in the blanks.
  \begin{multicols}{2}
    \begin{enumerate}[itemsep=0.42cm]
      \item 3.2 km = \blank\ m
      \item 60 cm = \blank\ m
      \item 1.5 m = \blank\ cm
      \item 9 mm = \blank\ cm
      \item 0.4 m = \blank\ mm
      \item 750 m = \blank\ km
    \end{enumerate}
  \end{multicols}

\item Squares and square roots.
  \begin{multicols}{4}
    \begin{enumerate}[itemsep=0.5cm]
      \item $4^2 = $ \blank[1.1cm]
      \item $9^2 = $ \blank[1.1cm]
      \item $11^2 = $ \blank[1.1cm]
      \item $30^2 = $ \blank[1.1cm]
      \item $0.2^2 = $ \blank[1.1cm]
      \item $1.5^2 = $ \blank[1.1cm]
      \item $\sqrt{49} = $ \blank[1.1cm]
      \item $\sqrt{400} = $ \blank[1.1cm]
    \end{enumerate}
  \end{multicols}

\item Use 1 m/s = 3.6 km/h.
  \begin{multicols}{2}
    \begin{enumerate}[itemsep=0.42cm]
      \item 5 m/s = \blank[1.8cm] km/h
      \item 15 m/s = \blank[1.8cm] km/h
      \item 72 km/h = \blank[1.8cm] m/s
      \item 90 km/h = \blank[1.8cm] m/s
    \end{enumerate}
  \end{multicols}

\item Distance = speed $\times$ time.
  \begin{enumerate}[itemsep=0.42cm]
    \item A student walks at 1.5 m/s for 20 s. How far? \hfill \blank
    \item A cyclist rides at 6 m/s for one minute. How far? \hfill \blank
    \item A car drives at 90 km/h for 20 minutes. How far? \hfill \blank
    \item How long does it take to run 100 m at 5 m/s? \hfill \blank
  \end{enumerate}

\item A runner jogs one lap of a 400 m track in 80 s, finishing where the lap began.

  \vspace{0.3cm}
  Average speed = \blank \hfill Average velocity = \blank

\newpage

\item Write each fraction as a decimal.
  \begin{multicols}{4}
    \begin{enumerate}[itemsep=0.5cm]
      \item $\dfrac{3}{5} = $ \blank[1.6cm]
      \item $\dfrac{9}{10} = $ \blank[1.6cm]
      \item $\dfrac{3}{20} = $ \blank[1.6cm]
      \item $\dfrac{1}{50} = $ \blank[1.6cm]
      \item $\dfrac{3}{25} = $ \blank[1.6cm]
      \item $\dfrac{5}{4} = $ \blank[1.6cm]
      \item $\dfrac{1}{9} = $ \blank[1.6cm]
      \item $\dfrac{4}{3} = $ \blank[1.6cm]
    \end{enumerate}
  \end{multicols}

\item Fill in the blanks.
  \begin{multicols}{2}
    \begin{enumerate}[itemsep=0.42cm]
      \item 0.5 hour = \blank[1.8cm] min
      \item 0.4 hour = \blank[1.8cm] min
      \item 2.5 min = \blank[1.8cm] s
      \item 90 s = \blank[1.8cm] min
      \item 6 min = \blank[1.8cm] h (decimal)
      \item 40 min = \blank[1.8cm] h (fraction)
      \item 1 week = \blank[1.8cm] hours
      \item 2.25 hours = \blank[1.2cm] h \blank[1.2cm] min
    \end{enumerate}
  \end{multicols}

\item Acceleration = change in velocity $\div$ time.
  \begin{enumerate}[itemsep=0.42cm]
    \item A car goes from 0 to 20 m/s in 5 s. Find its acceleration. \hfill \blank
    \item A bike speeds up from 2 m/s to 8 m/s in 3 s. Find its acceleration.
      \hfill \blank
    \item A car at 20 m/s brakes, losing 5 m/s every second. How many seconds
      does it take to stop? \hfill \blank
    \item For the car in (a): average velocity = \blank[1.8cm] \hfill distance in
      the 5 s = \blank[1.8cm]
  \end{enumerate}

\item Use $g \approx 10$ m/s$^2$: a falling object gains 10 m/s of speed every second.
  \begin{enumerate}[itemsep=0.42cm]
    \item A stone is dropped. How fast is it going after 3 s? \hfill \blank
    \item How far has it fallen after 2 s? \; ($\frac{1}{2} \times 10 \times t^2$)
      \hfill \blank
    \item A ball is thrown straight up at 20 m/s. How many seconds does it take to
      reach the top? \hfill \blank
    \item At the top: velocity = \blank[1.8cm] \hfill acceleration =
      \blank[1.8cm]
  \end{enumerate}

\item A marathon is about 42 km, and 1 km $\approx \frac{5}{8}$ mile. About how many miles?
  \hfill \blank

\end{enumerate}
\end{document}
